% Lecture 6 notes: \input by notes.tex, \subimport'ed by the semester book.
% Exam questions (экзамен): 12 (разность квадратов, трёхшаговая процедура), 13 (B-гладкие
% числа, квадратичное решето), 14 (сложность и выбор параметров квадратичного решета, с выводом).
\begin{lecture}{6}{Разность квадратов и квадратичное решето}

\section{Зачем раскладывать на множители}

Стойкость RSA (лекция~5) держится на том, что по модулю $N = pq$ трудно найти $p$ и $q$: зная
их, противник вычисляет $\varphi(N) = (p - 1)(q - 1)$ и секретную экспоненту
$d = e^{-1} \bmod \varphi(N)$. Поэтому размер ключа RSA определяется тем, насколько быстро
работают \emph{лучшие известные} алгоритмы факторизации.

Алгоритмы разложения делятся на два класса.
\begin{itemize}
  \item \emph{Специальные} алгоритмы используют свойства множителей: пробное деление и
    $\rho$-метод Полларда быстры, если $N$ имеет \emph{маленький} простой делитель; $(p - 1)$-метод
    Полларда (лекция~5) "--- если $p - 1$ раскладывается на маленькие простые; метод Ферма (этот
    раздел) "--- если $p$ и $q$ \emph{близки}. Их существование объясняет требования к
    простым числам RSA.
  \item \emph{Универсальные} алгоритмы (метод Диксона, квадратичное решето, решето числового поля)
    имеют сложность, зависящую только от размера $N$. Все они основаны на одной идее "---
    \emph{разности квадратов}: найти $x, y$ с $x^2 \equiv y^2 \pmod N$, $x \not\equiv \pm y$.
\end{itemize}

Идея восходит к П.~Ферма (1643); М.~Крайчик (1920-е) заменил равенство $N = x^2 - y^2$
сравнением; М.~Моррисон и Дж.~Бриллхарт (1975) добавили линейную алгебру над $\GF{2}$;
Дж.~Диксон (1981) доказал субэкспоненциальную оценку, а К.~Померанс (1981) предложил
\emph{квадратичное решето}~\cite{pomerance-qs}, которое до появления решета числового поля
было самым быстрым алгоритмом факторизации. Подробная история "--- в~\cite{pomerance-two-sieves}.

Всюду ниже $N$ "--- нечётное составное число, не являющееся степенью простого (чётность и
точные степени проверяются быстро: $N = a^k$ проверяется извлечением корней степеней
$k \le \log_2 N$).

\section{Метод Ферма}

\begin{lemma}\label{l6:fermat-bij}
Пусть $N$ нечётно. Отображение
\[
  (a, b) \mapsto (x, y) = \Bigl(\frac{a + b}{2}, \frac{a - b}{2}\Bigr)
\]
"--- биекция между разложениями $N = ab$ с $a \ge b \ge 1$ и представлениями
$N = x^2 - y^2$ с целыми $x > 0$, $y \ge 0$. Обратное отображение: $a = x + y$, $b = x - y$.
\end{lemma}

\begin{proof}
Так как $N$ нечётно, $a$ и $b$ нечётны, поэтому $a \pm b$ чётны и $x, y$ "--- целые;
$x^2 - y^2 = \frac{(a + b)^2 - (a - b)^2}{4} = ab = N$. Обратно, если $N = x^2 - y^2 = (x - y)(x + y)$
с $x > 0$, $y \ge 0$, то $x + y \ge 1$, значит, и $x - y = N/(x + y) \ge 1$, и пара
$(a, b) = (x + y, x - y)$ "--- разложение с $a \ge b \ge 1$. Отображения взаимно обратны.
\end{proof}

Разложению $N = N \cdot 1$ соответствует тривиальное представление
$x = \frac{N + 1}{2}$, $y = \frac{N - 1}{2}$. Метод Ферма ищет \emph{нетривиальное}
представление, перебирая $x$ от $\ceil{\sqrt N}$ вверх.

\begin{algorithm}[метод Ферма]
\begin{algorithmic}[1]
\Require нечётное $N > 1$.
\Ensure разложение $N = (x - y)(x + y)$.
\State $x \gets \ceil{\sqrt N}$, \quad $r \gets x^2 - N$
\While{$r$ "--- не полный квадрат}
  \State $r \gets r + 2x + 1$, \quad $x \gets x + 1$ \Comment{$r = x^2 - N$}
\EndWhile
\State $y \gets \sqrt r$
\State \Return $(x - y,\ x + y)$
\end{algorithmic}
\end{algorithm}

\begin{theorem}\label{l6:fermat-steps}
Пусть $N = ab$, где $b$ "--- наибольший делитель $N$, не превосходящий $\sqrt N$, $a = N/b$.
Метод Ферма останавливается при $x = \frac{a + b}{2}$ и возвращает разложение $N = b \cdot a$,
выполнив
\[
  s = \frac{a + b}{2} - \ceil{\sqrt N} + 1
\]
проверок «$x^2 - N$ "--- квадрат?». При этом
\[
  \frac{a + b}{2} - \sqrt N = \frac{(\sqrt a - \sqrt b)^2}{2} = \frac{(a - b)^2}{2(\sqrt a + \sqrt b)^2}
  \le \frac{(a - b)^2}{8\sqrt N}.
\]
\end{theorem}

\begin{proof}
По лемме~\ref{l6:fermat-bij} значения $x$, при которых $x^2 - N$ "--- полный квадрат, "---
это в точности числа $\frac{a' + b'}{2}$ по всем разложениям $N = a'b'$, $a' \ge b'$. Так как
$(a' + b')^2 = (a' - b')^2 + 4N$, число $x = \frac{a' + b'}{2}$ возрастает вместе с разностью
$a' - b'$, т.\,е. убывает с ростом $b' \le \sqrt N$. Наименьшее из них соответствует наибольшему
делителю $b' = b \le \sqrt N$, и $\frac{a + b}{2} \ge \sqrt{ab} = \sqrt N$ (неравенство о средних),
т.\,е. $\frac{a+b}{2} \ge \ceil{\sqrt N}$. Алгоритм проверяет $x = \ceil{\sqrt N}, \ceil{\sqrt N} + 1, \dots$
по порядку, поэтому первым найдёт именно $x = \frac{a+b}{2}$, сделав $s$ проверок.

Равенство $\frac{a + b}{2} - \sqrt{ab} = \frac{(\sqrt a - \sqrt b)^2}{2}$ проверяется раскрытием
скобок; далее $\sqrt a - \sqrt b = \frac{a - b}{\sqrt a + \sqrt b}$ и
$(\sqrt a + \sqrt b)^2 = a + b + 2\sqrt N \ge 4\sqrt N$.
\end{proof}

\begin{corollary}\label{l6:fermat-close}
Если $N = pq$, $p < q$ "--- простые и $q - p \le C N^{1/4}$, то метод Ферма раскладывает $N$
не более чем за $C^2/8 + 1$ шагов. Если же $p$ мало ($p \le N^{1/2 - \varepsilon}$), то число
шагов $s > \frac{p + q}{2} - \sqrt N > \frac{N}{2p} - \sqrt N \ge \frac12 N^{1/2 + \varepsilon} - N^{1/2}$
"--- метод медленнее пробного деления.
\end{corollary}

Один шаг стоит одного сложения $\bigO(n)$-битных чисел ($n = \log_2 N$) и проверки на квадрат.
Последнюю почти всегда можно сделать без извлечения корня: квадраты по модулю $16$ "--- только
$0, 1, 4, 9$, по модулю $9$ "--- $0, 1, 4, 7$ и т.\,д.; по нескольким таким модулям отсеивается
подавляющее большинство кандидатов, а корень вычисляется лишь для оставшихся.

\begin{example}\label{l6:ex-fermat}
$N = 8051$: $\ceil{\sqrt{8051}} = 90$, $90^2 - 8051 = 49 = 7^2$, откуда
$8051 = (90 - 7)(90 + 7) = 83 \cdot 97$ "--- один шаг.
$N = 5959$: $\ceil{\sqrt{5959}} = 78$;
\[
  78^2 - N = 125, \qquad 79^2 - N = 282, \qquad 80^2 - N = 441 = 21^2,
\]
и $5959 = 59 \cdot 101$, $s = 3 = \frac{59 + 101}{2} - 78 + 1$.
\end{example}

\taskexample{fermat-trace}

\paragraph{Множители $kN$.} Если $q/p$ близко к рациональному числу $v/u$ с маленькими $u, v$,
то у числа $kN = (vp)(uq)$, $k = uv$, множители $vp$ и $uq$ близки, и метод Ферма, применённый к
$kN$, быстро находит $kN = x^2 - y^2$; делитель $N$ "--- это $\gcd(x - y, N)$. Например,
$N = 31007 = 101 \cdot 307$: для самого $N$ нужно $s = 204 - 177 + 1 = 28$ шагов, а для
$3N = 93021$ уже первый шаг даёт $305^2 - 93021 = 4 = 2^2$, $3N = 303 \cdot 307$ и
$\gcd(303, 31007) = 101$. Перебирая все $k \le N^{1/3}$, Р.~Леман (1974) получил
детерминированный алгоритм сложности $\bigO(N^{1/3})$ (без доказательства,
см.~\cite[разд.~5.1]{crandall-pomerance}).

\paragraph{Вывод для RSA.} Простые $p, q$ ключа RSA нельзя выбирать слишком близкими:
стандарт FIPS~186 требует $\abs{p - q} > 2^{n/2 - 100}$ для $n$-битного модуля. Это не
теоретическая предосторожность: в 2022~г. Х.~Бёк обнаружил принтеры, генерировавшие ключи RSA с
близкими $p$ и $q$ (оба простых искались рядом с одним и тем же случайным числом), "--- такие ключи раскладываются
методом Ферма за секунды.

\checkyourself{Почему число $N \equiv 2 \pmod 4$ нельзя представить в виде разности
квадратов? Сколько шагов сделает метод Ферма для простого $N$ и для $N = p^2$?}

\checkyourself{Метод Ферма находит разложение $N = ab$ с $b$, ближайшим к $\sqrt N$. Каким будет
результат для $N = 3 \cdot 5 \cdot 7 \cdot 11 = 1155$? Можно ли получить с его помощью полное
разложение на простые?}

\section{Сравнение квадратов}

Равенство $N = x^2 - y^2$ слишком требовательно. Достаточно сравнения.

\begin{lemma}\label{l6:cong}
Если $\congr{x^2}{y^2}{N}$ и $x \not\equiv \pm y \pmod N$, то $d = \gcd(x - y, N)$ "---
нетривиальный делитель $N$: $1 < d < N$.
\end{lemma}

\begin{proof}
$N \mid (x - y)(x + y)$. Если $d = N$, то $N \mid x - y$, т.\,е. $x \equiv y$. Если $d = 1$, то
из $N \mid (x - y)(x + y)$ следует $N \mid x + y$, т.\,е. $x \equiv -y$. Обе возможности
исключены условием.
\end{proof}

Насколько часто сравнение квадратов «полезно»? Ответ даёт подсчёт квадратных корней.

\begin{lemma}\label{l6:roots-pe}
Пусть $p$ "--- нечётное простое, $e \ge 1$, $p \nmid y$. Сравнение $\congr{x^2}{y^2}{p^e}$ имеет
ровно два решения: $x \equiv \pm y \pmod{p^e}$, и они различны.
\end{lemma}

\begin{proof}
$p^e \mid (x - y)(x + y)$. Если бы $p$ делило оба множителя, то $p$ делило бы их разность $2y$,
что невозможно ($p$ нечётно, $p \nmid y$). Значит, $p^e$ целиком делит один из множителей:
$x \equiv y$ или $x \equiv -y$. Наконец, $y \not\equiv -y \pmod{p^e}$, так как иначе $p \mid 2y$.
\end{proof}

\begin{theorem}\label{l6:half}
Пусть $N = p_1^{e_1} \cdots p_k^{e_k}$ нечётно, $k \ge 2$, и $\gcd(y, N) = 1$. Тогда
сравнение $\congr{x^2}{y^2}{N}$ имеет ровно $2^k$ решений по модулю $N$, и только два из них
"--- $x \equiv \pm y$. Следовательно, если $x$ выбран равновероятно среди этих решений, то
$\gcd(x - y, N)$ нетривиален с вероятностью $1 - 2^{1-k} \ge \frac12$.
\end{theorem}

\begin{proof}
По китайской теореме об остатках $\congr{x^2}{y^2}{N}$ равносильно системе
$\congr{x^2}{y^2}{p_i^{e_i}}$, $i = 1, \dots, k$. По лемме~\ref{l6:roots-pe} каждое сравнение
имеет два решения $x \equiv \varepsilon_i y \pmod{p_i^{e_i}}$, $\varepsilon_i = \pm 1$, и каждому
набору знаков $(\varepsilon_1, \dots, \varepsilon_k)$ по КТО отвечает ровно одно решение по модулю
$N$; всего $2^k$ решений. Решения $x \equiv y$ и $x \equiv -y$ "--- это наборы $(1, \dots, 1)$ и
$(-1, \dots, -1)$; для остальных $2^k - 2$ наборов $x \not\equiv \pm y$, и по
лемме~\ref{l6:cong} $\gcd(x - y, N)$ нетривиален. Доля таких решений равна
$1 - 2/2^k \ge 1 - 2/4 = \frac12$.
\end{proof}

\begin{example}
$N = 21$, $y = 2$: решения $\congr{x^2}{4}{21}$ "--- $x \in \{2, 5, 16, 19\}$. Для $x = 5$:
$\gcd(5 - 2, 21) = 3$; для $x = 16 \equiv -5$: $\gcd(14, 21) = 7$; $x = 2, 19$ дают $21$ и $1$.
\end{example}

Для $N = p^e$ (одно простое) теорема ничего не даёт: корней только два. Поэтому степени
простых исключают заранее. В алгоритмах ниже $x$ вычисляется детерминированно, и утверждение
«с вероятностью $\ge \frac12$» "--- \emph{эвристика}: предполагается, что найденное $x$ ведёт себя
как случайный корень из $y^2$. Основание для неё такое: алгоритм использует числа $a_i$
только через их квадраты $a_i^2 \bmod N$, а каждый квадрат имеет $2^k$ равноправных корней.

\checkyourself{Почему для $N = pqr$ (три различных простых) вероятность успеха в
теореме~\ref{l6:half} равна $3/4$? Что произойдёт, если $\gcd(y, N) > 1$?}

\section{Трёхшаговая процедура}

Как находить сравнения квадратов? Пусть даны $a_i$, для которых вычеты
$c_i \equiv a_i^2 \pmod N$ раскладываются на маленькие простые. Ни одно $c_i$ само по себе не
квадрат, но произведение подходящих $c_i$ может им быть, и тогда
\[
  X = \prod_{i \in S} a_i, \qquad Y = \sqrt{\prod_{i \in S} c_i}, \qquad
  X^2 = \prod_{i \in S} a_i^2 \equiv \prod_{i \in S} c_i = Y^2 \pmod N.
\]

\begin{algorithm}[трёхшаговая процедура разложения, \cite{hps}]
\begin{algorithmic}[1]
\Require нечётное составное $N$, не степень простого.
\State \textbf{Построение соотношений.} Найти $a_1, \dots, a_r$, для которых
  $c_i \equiv a_i^2 \pmod N$ раскладываются на простые из небольшого множества
  $\{p_1, \dots, p_t\}$ (\emph{факторной базы}): $c_i = \prod_j p_j^{e_{ij}}$.
\State \textbf{Исключение.} Найти непустое $S \subseteq \{1, \dots, r\}$, для которого
  $\prod_{i \in S} c_i$ "--- полный квадрат, т.\,е. все суммы $\sum_{i \in S} e_{ij}$ чётны.
\State \textbf{НОД.} $X \gets \prod_{i \in S} a_i \bmod N$,
  $Y \gets \prod_j p_j^{\frac12\sum_{i \in S} e_{ij}} \bmod N$, $d \gets \gcd(X - Y, N)$.
  Если $1 < d < N$ "--- вернуть $d$, иначе взять другое $S$.
\end{algorithmic}
\end{algorithm}

\begin{example}[\cite{pomerance-two-sieves}]\label{l6:ex-1649}
$N = 1649$, $\ceil{\sqrt N} = 41$:
\[
  41^2 \equiv 32 = 2^5, \qquad 42^2 \equiv 115 = 5 \cdot 23, \qquad 43^2 \equiv 200 = 2^3 \cdot 5^2
  \pmod{1649}.
\]
Произведение $32 \cdot 200 = 6400 = 80^2$ "--- квадрат, поэтому $X = 41 \cdot 43 = 1763 \equiv 114$,
$Y = 80$, $\congr{114^2}{80^2}{1649}$ и $\gcd(114 - 80, 1649) = \gcd(34, 1649) = 17$:
$1649 = 17 \cdot 97$.
\end{example}

\subsection{Линейная алгебра над $\GF{2}$}

\begin{definition}
\emph{Вектор показателей} числа $c = \prod_{j=1}^{t} p_j^{e_j}$ по модулю $2$ "--- это
$v(c) = (e_1 \bmod 2, \dots, e_t \bmod 2) \in \GF{2}^{t}$.
\end{definition}

Ясно, что $v(c c') = v(c) + v(c')$ и что $c$ "--- полный квадрат тогда и только тогда, когда
$v(c) = 0$. Поэтому шаг «исключение» "--- поиск линейной зависимости
$\sum_{i \in S} v(c_i) = 0$ между векторами $v(c_1), \dots, v(c_r)$ над полем $\GF{2}$.

\begin{proposition}\label{l6:deps}
Если $r > t$, то среди любых $r$ векторов $v_1, \dots, v_r \in \GF{2}^t$ найдётся непустая
зависимость; более того, непустых множеств $S$ с $\sum_{i \in S} v_i = 0$ не меньше
$2^{r - t} - 1$.
\end{proposition}

\begin{proof}
Множества $S$ с $\sum_{i \in S} v_i = 0$ "--- это векторы ядра линейного отображения
$\GF{2}^r \to \GF{2}^t$, $(s_1, \dots, s_r) \mapsto \sum s_i v_i$. По теореме о ранге
$\dim \ker = r - \mathrm{rank} \ge r - t$, т.\,е. в ядре не меньше $2^{r-t}$ векторов, из них
ненулевых "--- не меньше $2^{r-t} - 1$.
\end{proof}

Если каждая зависимость независимо даёт делитель с вероятностью $\ge \frac12$
(теорема~\ref{l6:half}), то при $r = t + 10$ соотношениях все $\ge 10$ базисных зависимостей
ядра окажутся неудачными с вероятностью $\le 2^{-10}$. Поэтому на практике собирают
\emph{чуть больше}, чем $t$, соотношений.

Зависимости находят методом Гаусса: строки $v_i$ приводят к ступенчатому виду, запоминая, из
каких исходных строк составлена каждая текущая строка; строки, ставшие нулевыми, и дают
множества $S$. Над $\GF{2}$ сложение строк "--- это XOR, и строку из $t$ битов можно хранить в
$\ceil{t/64}$ машинных словах; сложность "--- $\bigO(t^3)$ битовых операций, или $\bigO(t^3/64)$
словных. Матрицы соотношений \emph{разреженные}: в строке не больше $\log_2 N$ единиц. Для них
есть методы сложности $\bigO(t^2 \cdot w)$, где $w$ "--- среднее число единиц в строке:
блочный метод Ланцоша (П.~Монтгомери, 1995) и блочный метод Видемана; перед ними матрицу
сжимают «структурированным исключением». Именно они используются в рекордных вычислениях
(без доказательства, см.~\cite[разд.~6.1.3]{crandall-pomerance}).

\taskexample{squares-combination}

\checkyourself{В примере~\ref{l6:ex-1649} зависимость нашлась «на глаз». Запишите векторы
показателей для $41^2, 42^2, 43^2$ по базе $\{2, 5, 23\}$ и найдите её методом Гаусса.}

\section{Гладкие числа}

Шаг «построение соотношений» требует чисел, раскладывающихся на маленькие простые.

\begin{definition}
Число $n$ называется \emph{$B$-гладким}, если все его простые делители не превосходят $B$.
Через $\psi(X, B)$ обозначается количество $B$-гладких чисел $1 \le n \le X$.
\end{definition}

Например, $\psi(100, 5) = 34$ (числа вида $2^a 3^b 5^c \le 100$, включая $1$). Гладкие числа
встречаются тем реже, чем меньше $B$ по сравнению с $X$; естественный параметр "--- отношение
\[
  u = \frac{\ln X}{\ln B}, \qquad \text{т.\,е. } B = X^{1/u}.
\]

\begin{proposition}\label{l6:rho2}
Для $1 \le u \le 2$ при $X \to \infty$
\[
  \frac{\psi(X, X^{1/u})}{X} \to 1 - \ln u.
\]
\end{proposition}

\begin{proof}
Пусть $B = X^{1/u} \ge \sqrt X$. Число $n \le X$ не является $B$-гладким, если у него есть
простой делитель $p > B$; такой делитель единственный, так как $p^2 > B^2 \ge X$. Поэтому
\[
  X - \psi(X, B) = \sum_{B < p \le X} \floor*{\frac{X}{p}} = X \sum_{B < p \le X} \frac 1p
    + \bigO(\pi(X)).
\]
По теореме Мертенса $\sum_{p \le x} \frac1p = \ln\ln x + M + o(1)$ (без доказательства,
см.~\cite[разд.~22.7]{hardy-wright}), поэтому
$\sum_{B < p \le X} \frac1p = \ln\ln X - \ln\ln X^{1/u} + o(1) = \ln u + o(1)$, а
$\pi(X) = o(X)$. Делим на $X$.
\end{proof}

Например, доля $1000$-гладких чисел среди $n \le 10^6$ ($u = 2$) равна $0{,}344$, а предел
$1 - \ln 2 \approx 0{,}307$. Предельная функция $\rho(u)$ (\emph{функция Дикмана}) при всех $u$
определяется уравнением $u \rho'(u) = -\rho(u - 1)$ и убывает примерно как $u^{-u}$. Для
алгоритмов нужен случай, когда $u$ растёт вместе с $X$.

\begin{theorem}[Кэнфилд"--~Эрдёш"--~Померанс~\cite{canfield-erdos-pomerance}, без доказательства]
\label{l6:cep}
Для любого $\varepsilon > 0$ равномерно по $3 \le u \le (1 - \varepsilon)\frac{\ln X}{\ln\ln X}$
\[
  \psi(X, X^{1/u}) = X \cdot u^{-u(1 + o(1))}, \qquad X \to \infty.
\]
\end{theorem}

Иначе говоря, «случайное» число $n \le X$ является $X^{1/u}$-гладким с вероятностью
$u^{-u(1+o(1))}$. Удобно записывать такие величины через \emph{субэкспоненциальную функцию}
\[
  L(X) = e^{\sqrt{\ln X \ln\ln X}}, \qquad
  L_X[\alpha, c] = e^{(c + o(1)) (\ln X)^{\alpha} (\ln\ln X)^{1 - \alpha}}.
\]
При $\alpha = 0$ получаем $(\ln X)^{c + o(1)}$ "--- полиномиальную сложность от длины $X$, при
$\alpha = 1$ "--- $X^{c + o(1)}$, экспоненциальную; $L(X)^c = L_X[\frac12, c]$ лежит между ними.
Заметим, что $L(X)^c = X^{o(1)}$ при любом фиксированном $c$.

\begin{corollary}\label{l6:psiL}
Для любых фиксированных $\beta, \gamma > 0$ при $X = N^{\gamma + o(1)}$ и $B = L(N)^{\beta}$
\[
  \frac{\psi(X, B)}{X} = L(N)^{-\frac{\gamma}{2\beta} + o(1)}.
\]
В частности, $\psi(X, L(X)^{\beta}) = X \cdot L(X)^{-\frac{1}{2\beta} + o(1)}$.
\end{corollary}

\begin{proof}
Пусть $\ell = \ln N$. Тогда $\ln X = (\gamma + o(1))\ell$, $\ln B = \beta\sqrt{\ell \ln \ell}$ и
\[
  u = \frac{\ln X}{\ln B} = \Bigl(\frac{\gamma}{\beta} + o(1)\Bigr)\sqrt{\frac{\ell}{\ln\ell}}, \qquad
  \ln u = \frac12 \ln\ell - \frac12\ln\ln\ell + \bigO(1) = \Bigl(\frac12 + o(1)\Bigr)\ln \ell.
\]
Условие теоремы~\ref{l6:cep} выполнено: $u \asymp \sqrt{\ell / \ln \ell} = o(\ell/\ln\ell)$.
Поэтому
\[
  u^{-u(1 + o(1))} = e^{-(1 + o(1)) u \ln u}
  = e^{-\left(\frac{\gamma}{2\beta} + o(1)\right)\sqrt{\ell \ln\ell}} = L(N)^{-\frac{\gamma}{2\beta} + o(1)}.
\]
Второе утверждение "--- случай $N = X$, $\gamma = 1$.
\end{proof}

\checkyourself{Чему равны $\psi(X, 2)$ и $\psi(X, 3)$ с точностью до $\bigO(1)$ и
$\bigO(\ln X)$? Почему их рост полиномиален по $\ln X$, а не по $X$?}

\section{Метод Диксона}

Простейший способ построения соотношений "--- брать $a$ случайно.

\begin{algorithm}[Диксон]\label{l6:dixon-alg}
\begin{algorithmic}[1]
\Require $N$, граница гладкости $B$; факторная база "--- все простые $p_1 < \cdots < p_t \le B$.
\While{собрано меньше $t + 10$ соотношений}
  \State выбрать случайное $a \in [1, N - 1]$, $c \gets a^2 \bmod N$
  \State если $c$ "--- $B$-гладкое (пробное деление на $p_1, \dots, p_t$), сохранить $(a, c)$
\EndWhile
\State исключение над $\GF{2}$ и НОД, как в трёхшаговой процедуре
\end{algorithmic}
\end{algorithm}

\begin{theorem}[эвристическая оценка]\label{l6:dixon}
Пусть числа $c = a^2 \bmod N$ ведут себя как случайные числа из $[1, N]$. Тогда при
$B = L(N)^{\beta}$ метод Диксона выполняет $L(N)^{2\beta + \frac{1}{2\beta} + o(1)}$ операций,
и оптимум $\beta = \frac12$ даёт $L(N)^{2 + o(1)}$. Если проверка гладкости стоит $L(N)^{o(1)}$
операций, то сложность равна $L(N)^{\beta + \frac{1}{2\beta} + o(1)}$, оптимум достигается при
$B = L(N)^{1/\sqrt2}$ и равен $L(N)^{\sqrt 2 + o(1)}$.
\end{theorem}

\begin{proof}[Вывод]
По следствию~\ref{l6:psiL} ($\gamma = 1$) число $c$ гладко с вероятностью
$P = L(N)^{-\frac{1}{2\beta} + o(1)}$. Нужно $t + 10 \le \pi(B) + 10 \le B$ соотношений, т.\,е.
в среднем $B / P = L(N)^{\beta + \frac{1}{2\beta} + o(1)}$ попыток. Пробное деление в одной
попытке стоит $\pi(B) = L(N)^{\beta + o(1)}$ делений, отсюда первая оценка. Линейная алгебра
(разреженная) стоит $t^{2 + o(1)} = L(N)^{2\beta + o(1)}$ и не превосходит затрат на поиск.
Функция $f(\beta) = 2\beta + \frac{1}{2\beta}$ минимальна при $f'(\beta) = 2 - \frac{1}{2\beta^2} = 0$,
$\beta = \frac12$, $f = 2$; функция $g(\beta) = \beta + \frac{1}{2\beta}$ "--- при
$1 - \frac{1}{2\beta^2} = 0$, $\beta = \frac{1}{\sqrt2}$, $g = \sqrt2$ (и $2\beta = \sqrt 2$ "---
линейная алгебра не хуже).
\end{proof}

Диксон~\cite{dixon} доказал оценку вида $L(N)^{c}$ \emph{строго}, без эвристик "--- это первый
алгоритм факторизации с доказанной субэкспоненциальной сложностью. Его недостаток: числа $c$ имеют
размер $N$, и гладкие среди них редки.

\section{Квадратичное решето}

Квадратичное решето (QS) улучшает метод Диксона двумя приёмами.
\begin{enumerate}
  \item \emph{Малые вычеты}: $a$ берут рядом с $\sqrt N$, тогда $c = a^2 - N$ имеет размер
    порядка $\sqrt N$, а не $N$, и гладким оказывается гораздо чаще.
  \item \emph{Просеивание}: значения $c$ "--- значения многочлена, и делимость на $p$ зависит
    только от остатка аргумента по модулю $p$. Это позволяет заменить пробное деление
    каждого числа решетом, как в решете Эратосфена.
\end{enumerate}

Пусть $m = \ceil{\sqrt N}$ и
\[
  F(t) = (t + m)^2 - N, \qquad t = 0, 1, 2, \dots
\]
Тогда $\congr{(t + m)^2}{F(t)}{N}$, и при $0 \le t \le M$
\[
  0 \le F(t) < (\sqrt N + t + 1)^2 - N = 2(t + 1)\sqrt N + (t + 1)^2 \le 2(M + 1)\sqrt N + (M + 1)^2.
\]

\subsection{Факторная база}

\begin{proposition}\label{l6:fb}
Пусть $p$ "--- нечётное простое, $p \nmid N$.
\begin{enumerate}
  \item Если $p \mid F(t)$ хотя бы при одном $t$, то $\legendre{N}{p} = 1$.
  \item Если $\legendre{N}{p} = 1$ и $\congr{r^2}{N}{p}$, то $p \mid F(t)$ тогда и только тогда,
    когда $t \equiv r - m$ или $t \equiv -r - m \pmod p$ "--- это два различных класса вычетов.
\end{enumerate}
\end{proposition}

\begin{proof}
1. $p \mid F(t)$ означает $\congr{(t + m)^2}{N}{p}$, т.\,е. $N$ "--- квадрат по модулю $p$,
причём $N \not\equiv 0$; значит, $\legendre{N}{p} = 1$.
2. $p \mid F(t)$ $\iff$ $\congr{(t + m)^2}{r^2}{p}$ $\iff$ $t + m \equiv \pm r \pmod p$
(лемма~\ref{l6:roots-pe} при $e = 1$); $r \not\equiv -r$, так как $p \nmid 2r$.
\end{proof}

Таким образом, нечётные простые с $\legendre{N}{p} = -1$ \emph{никогда} не делят значения $F$,
и включать их в факторную базу бессмысленно. \emph{Факторная база} квадратичного решета:
\[
  \mathcal{F} = \{2\} \cup \Bigl\{p \le B \text{ "--- нечётное простое} : \legendre{N}{p} = 1\Bigr\}
\]
(простое $p \mid N$ сразу даёт делитель). Символ Лежандра вычисляется по критерию Эйлера или
через символ Якоби; доля простых с $\legendre{N}{p} = 1$ асимптотически равна $\frac12$, так
что $\abs{\mathcal F} \approx \frac12\pi(B)$. Корни $\pm r$ из $N$ по модулю $p$ находят методами
Тонелли"--~Шенкса или Чиполлы (лекция~7), а по модулю $p^e$ "--- подъёмом Гензеля (там же):
при $p \nmid 2N$ корней по модулю $p^e$ тоже ровно два (лемма~\ref{l6:roots-pe}). Для $p = 2$
(при нечётном $N$) $F(t)$ чётно ровно при нечётном $t + m$; классы по модулю $2^e$ находят
перебором.

\subsection{Просеивание}

\begin{algorithm}[просеивание, точный вариант]\label{l6:sieve}
\begin{algorithmic}[1]
\Require $N$, $m = \ceil{\sqrt N}$, длина интервала $M$, факторная база $\mathcal F$.
\Ensure все $B$-гладкие значения $F(t)$, $0 \le t < M$, с разложениями.
\State $A[t] \gets F(t)$ для $0 \le t < M$
\For{$p \in \mathcal F$}
  \For{$e = 1, 2, \dots$, пока $p^e \le \max_t F(t)$}
    \For{каждый корень $\rho$ сравнения $\congr{(\rho + m)^2}{N}{p^e}$, $0 \le \rho < p^e$}
      \For{$t = \rho, \rho + p^e, \rho + 2p^e, \dots < M$}
        \State $A[t] \gets A[t] / p$ \Comment{$p^e \mid F(t)$: делим на очередной множитель $p$}
      \EndFor
    \EndFor
  \EndFor
\EndFor
\State \Return $\{t : A[t] = 1\}$ \Comment{$F(t)$ "--- $B$-гладкие}
\end{algorithmic}
\end{algorithm}

Позиция $t$ делится на $p$ при обходе по модулю $p$, ещё раз "--- по модулю $p^2$ (если
$p^2 \mid F(t)$) и т.\,д., поэтому в итоге из $A[t]$ удалена ровно та степень $p$, которая
входит в $F(t)$, и $A[t] = 1$ тогда и только тогда, когда $F(t)$ раскладывается по базе.

На практике деление заменяют сложением \emph{логарифмов}: $A[t] \gets 0$; к $A[t]$ прибавляют
округлённое $\log_2 p$ для каждой степени $p^e \mid F(t)$, после чего кандидатами считаются
позиции, где $A[t] \ge \log_2 F(t) - T$ (порог $T$ учитывает округление и пропущенные малые
степени). Кандидатов немного, и их раскладывают пробным делением. Логарифмы занимают один
байт на позицию, и интервал обрабатывают кусками, помещающимися в кэш процессора.

\begin{lemma}\label{l6:sieve-cost}
Просеивание интервала длины $M$ по базе из простых $p \le B$ выполняет
$\bigO(M \ln\ln B + \pi(B) \log N)$ операций над элементами массива.
\end{lemma}

\begin{proof}
Для простого $p$ и степени $p^e$ обходится не больше двух классов (для $p = 2$ "--- не больше
четырёх), в каждом $\le M/p^e + 1$ позиций. Сумма по $e \ge 1$ даёт
$\bigO(M/p) + \bigO(\log_p N)$ на одно простое (степеней не больше $\log_2 N$). По теореме Мертенса
$\sum_{p \le B} \frac 1p = \ln\ln B + \bigO(1)$, откуда оценка.
\end{proof}

Итак, просеивание обходится в $\bigO(\ln\ln B)$ операций \emph{на одно значение} $F(t)$ "---
практически константу, тогда как пробное деление одного числа стоит $\pi(B)$ делений.

\begin{example}\label{l6:ex-qs}
$N = 9553$, $m = \ceil{\sqrt{9553}} = 98$ ($97^2 = 9409$, $98^2 = 9604$),
$F(t) = (t + 98)^2 - 9553$. Возьмём $B = 17$. Критерий Эйлера: $\legendre{9553}{p} = 1$ для
$p = 3, 11, 17$ и $-1$ для $p = 5, 7, 13$, так что $\mathcal F = \{2, 3, 11, 17\}$.
Позиции просеивания ($t \equiv \pm r - m \pmod p$):
\[
  p = 2: t \equiv 1; \qquad p = 3: t \equiv 0, 2; \qquad p = 11: t \equiv 5, 8; \qquad
  p = 17: t \equiv 0, 8.
\]
После деления на все степени простых из базы ($t = 0, \dots, 19$; «ост.» "--- то, что осталось
в $A[t]$):
\[
\small
\begin{array}{r|r|l|r}
t & F(t) & \text{из базы} & \text{ост.} \\ \hline
0 & 51 & 3 \cdot 17 & \mathbf{1} \\
1 & 248 & 2^3 & 31 \\
2 & 447 & 3 & 149 \\
3 & 648 & 2^3 \cdot 3^4 & \mathbf{1} \\
4 & 851 & & 851 \\
5 & 1056 & 2^5 \cdot 3 \cdot 11 & \mathbf{1} \\
6 & 1263 & 3 & 421 \\
7 & 1472 & 2^6 & 23 \\
8 & 1683 & 3^2 \cdot 11 \cdot 17 & \mathbf{1} \\
9 & 1896 & 2^3 \cdot 3 & 79
\end{array}
\qquad
\begin{array}{r|r|l|r}
t & F(t) & \text{из базы} & \text{ост.} \\ \hline
10 & 2111 & & 2111 \\
11 & 2328 & 2^3 \cdot 3 & 97 \\
12 & 2547 & 3^2 & 283 \\
13 & 2768 & 2^4 & 173 \\
14 & 2991 & 3 & 997 \\
15 & 3216 & 2^4 \cdot 3 & 67 \\
16 & 3443 & 11 & 313 \\
17 & 3672 & 2^3 \cdot 3^3 \cdot 17 & \mathbf{1} \\
18 & 3903 & 3 & 1301 \\
19 & 4136 & 2^3 \cdot 11 & 47
\end{array}
\]
Гладких значений пять; векторы показателей по модулю $2$ (порядок $2, 3, 11, 17$):
\[
\begin{aligned}
  v(F(0)) &= (0, 1, 0, 1), & v(F(3)) &= (1, 0, 0, 0), & v(F(5)) &= (1, 1, 1, 0),\\
  v(F(8)) &= (0, 0, 1, 1), & v(F(17)) &= (1, 1, 0, 1).
\end{aligned}
\]
Зависимость $v(F(0)) + v(F(3)) + v(F(17)) = 0$: $X = 98 \cdot 101 \cdot 115 \bmod N = 1463$,
$Y = 2^3 \cdot 3^4 \cdot 17 \bmod N = 11016 - 9553 = 1463$. Здесь $X \equiv Y$, и
$\gcd(X - Y, N) = N$ "--- неудача. Зависимость $v(F(5)) + v(F(8)) + v(F(17)) = 0$:
$X = 103 \cdot 106 \cdot 115 \bmod N = 4127$, $Y = 2^4 \cdot 3^3 \cdot 11 \cdot 17 \bmod N = 4360$,
$\gcd(4127 - 4360, 9553) = \gcd(233, 9553) = 233$, и $9553 = 41 \cdot 233$.
\end{example}

\taskexample{qs-small}

\subsection{Алгоритм целиком}

\begin{algorithm}[квадратичное решето]\label{l6:qs}
\begin{algorithmic}[1]
\Require нечётное составное $N$, не степень простого.
\State выбрать $B$ (см. ниже) и длину интервала $M$; $m \gets \ceil{\sqrt N}$
\State $\mathcal F \gets \{2\} \cup \{p \le B : \legendre{N}{p} = 1\}$; если некоторое $p \le B$
  делит $N$ "--- вернуть $p$
\For{$p \in \mathcal F$}
  \State найти корни $\pm r_p$ из $N$ по модулю $p$ (и $p^e$)
\EndFor
\State просеять $F(t) = (t + m)^2 - N$ на $[0, M)$ (алгоритм~\ref{l6:sieve}); если гладких
  значений меньше $\abs{\mathcal F} + 10$ "--- продолжить на следующем интервале
\State найти зависимости между векторами показателей над $\GF{2}$
\For{каждой зависимости $S$}
  \State $X \gets \prod_{t \in S} (t + m) \bmod N$, $Y \gets \sqrt{\prod_{t \in S} F(t)} \bmod N$
  \State если $1 < \gcd(X - Y, N) < N$ "--- вернуть этот делитель
\EndFor
\end{algorithmic}
\end{algorithm}

\paragraph{Улучшения.} 1) Просеивают по симметричному интервалу $-M \le t \le M$: значения $F(t)$
вдвое меньше по модулю, а знак учитывают, добавив в базу «простое» $p_0 = -1$. 2) \emph{Много
многочленов} (MPQS, П.~Монтгомери): вместо одного $F$ берут $F_{a,b}(t) = (at + b)^2 - N$ с
$a \approx \sqrt{2N}/M$, $b^2 \equiv N \pmod a$; тогда $F_{a,b}(t)/a$ остаются маленькими на
коротких интервалах, и интервал не приходится удлинять. SIQS ускоряет смену многочленов.
3) \emph{Большие простые}: значения вида (гладкое)$\cdot q$ с одним простым $q \in (B, B^2)$
сохраняют; два таких значения с одним $q$ дают полное соотношение.

\checkyourself{Почему в квадратичном решете не нужно проверять на гладкость числа $F(t)$ с
$t$, не попавшими ни в один класс $t \equiv \pm r_p - m$? Какое простое из базы всегда делит
$F(t)$ при нечётном $t + m$?}

\section{Сложность квадратичного решета и выбор параметров}

Строгой оценки для QS нет: она опирается на эвристическое предположение.

\begin{quote}
\textbf{(H)} Значения $F(t)$, $0 \le t \le M$, являются $B$-гладкими с той же вероятностью, что
и случайные числа того же размера.
\end{quote}

\begin{theorem}[эвристическая]\label{l6:qs-complexity}
При предположении (H) квадратичное решето с параметрами
\[
  B = L(N)^{1/2}, \qquad M = L(N)^{1 + o(1)}
\]
раскладывает $N$ за $L(N)^{1 + o(1)} = L_N[\frac12, 1]$ операций, используя
$L(N)^{1/2 + o(1)}$ памяти (не считая массива решета, который обрабатывают кусками). Это
значение $B$ оптимально: при $B = L(N)^{\beta}$ сложность равна
$L(N)^{\max(\beta + \frac{1}{4\beta},\ 2\beta) + o(1)}$.
\end{theorem}

\begin{proof}[Вывод]
Положим $B = L(N)^{\beta}$, $\beta > 0$, и будем искать $\beta$, минимизирующее время. Всюду
$\ell = \ln N$, а $L = L(N)$.

\emph{1. Размер значений.} При $M = L^{\mu}$ (какое-то фиксированное $\mu$) имеем
$F(t) < 2(M + 1)\sqrt N + (M + 1)^2 = N^{1/2} L^{\bigO(1)} = N^{1/2 + o(1)}$, так как $L^c = N^{o(1)}$.

\emph{2. Вероятность гладкости.} По (H) и следствию~\ref{l6:psiL} с $\gamma = \frac12$:
\[
  P = \frac{\psi(N^{1/2 + o(1)}, L^{\beta})}{N^{1/2 + o(1)}} = L^{-\frac{1}{4\beta} + o(1)}.
\]
Сравните с методом Диксона: там $\gamma = 1$ и $P = L^{-\frac{1}{2\beta}}$ "--- уменьшение
значений с $N$ до $\sqrt N$ «вдвое уменьшает» показатель.

\emph{3. Длина интервала.} Нужно $\abs{\mathcal F} + 10 \le \pi(B) + 10 = L^{\beta + o(1)}$
гладких значений, поэтому
\[
  M = \frac{L^{\beta + o(1)}}{P} = L^{\beta + \frac{1}{4\beta} + o(1)}.
\]

\emph{4. Время просеивания.} По лемме~\ref{l6:sieve-cost} это
$\bigO(M \ln\ln B + \pi(B)\ell) = L^{\beta + \frac{1}{4\beta} + o(1)}$
(множители $\ln\ln B$ и $\ell$ "--- это $L^{o(1)}$). Сюда же входит поиск корней $\pm r_p$ для
всех $p \le B$: $\pi(B) \cdot \ell^{\bigO(1)} = L^{\beta + o(1)}$, и пробное деление
$L^{\beta + o(1)}$ найденных гладких значений по $\pi(B)$ простым: $L^{2\beta + o(1)}$.

\emph{5. Линейная алгебра.} Матрица имеет размер $\approx \pi(B) \times \pi(B)$ и не больше
$\log_2 N$ единиц в строке. Разреженные методы (Ланцош, Видеман) требуют
$\pi(B)^{2} \ell^{\bigO(1)} = L^{2\beta + o(1)}$ операций, метод Гаусса "--- $L^{3\beta + o(1)}$.

\emph{6. Шаг НОД} стоит $\ell^{\bigO(1)}$ на зависимость и не влияет на оценку.

\emph{7. Оптимизация.} Итого (с разреженной алгеброй) $T(\beta) = L^{\max(f(\beta), 2\beta) + o(1)}$,
$f(\beta) = \beta + \frac{1}{4\beta}$. Функция $f$ выпукла,
$f'(\beta) = 1 - \frac{1}{4\beta^2} = 0$ при $\beta = \frac12$, и $f(\frac12) = 1$; при этом
$2\beta = 1$ тоже. При $\beta < \frac12$ растёт $f$, при $\beta > \frac12$ "--- $2\beta$, поэтому
$\beta = \frac12$ "--- точка минимума максимума, и $T = L^{1 + o(1)}$, $M = L^{f(1/2)} = L^{1 + o(1)}$.
Память "--- матрица из $L^{\beta + o(1)} = L^{1/2 + o(1)}$ строк по $\ell^{\bigO(1)}$ ненулевых
элементов.
\end{proof}

\begin{remark}
С плотным методом Гаусса оптимум другой: $\max(\beta + \frac{1}{4\beta}, 3\beta)$ минимален, когда
$\beta + \frac{1}{4\beta} = 3\beta$, т.\,е. $\beta = \frac{1}{2\sqrt2}$, и сложность равна
$L^{3/(2\sqrt2)} \approx L^{1{,}06}$. Именно поэтому разреженная линейная алгебра "--- обязательная
часть рекордных вычислений.
\end{remark}

\paragraph{Числа.} Для $N \approx 10^{100}$: $L(N) \approx 2{,}3 \cdot 10^{15}$,
$L(N)^{1/2} \approx 4{,}8 \cdot 10^{7}$. На практике $B$ берут в десятки раз меньше: слагаемые
$o(1)$ в показателе существенны, а MPQS и большие простые меняют баланс. В
примере~\ref{l6:ex-qs} $L(9553) \approx 90{,}5$ и $L^{1/2} \approx 9{,}5$; но при $B = 11$ база
$\{2, 3, 11\}$ даёт на том же интервале лишь два гладких значения, $F(3)$ и $F(5)$, с
независимыми векторами $(1, 0, 0)$ и $(1, 1, 1)$, "--- зависимостей нет. Поэтому мы взяли $B = 17$
(для таких маленьких $N$ асимптотический выбор $B$ "--- лишь ориентир).

\paragraph{Сравнение алгоритмов.} Для $N$ без специальной структуры (через $p$ обозначен
наименьший простой делитель $N$):
\[
\begin{array}{l|l}
\text{алгоритм} & \text{сложность} \\ \hline
\text{пробное деление} & \bigO(p) \le N^{1/2} \\
\rho\text{-метод Полларда} & \bigO(\sqrt p) \le N^{1/4} \text{ (эвристически)} \\
\text{метод Ферма} & \bigO\bigl((q - p)^2 / \sqrt N\bigr) \text{, в худшем случае } \bigO(N) \\
\text{метод Диксона} & L_N[\tfrac12, \sqrt 2] \\
\text{квадратичное решето} & L_N[\tfrac12, 1] \\
\text{метод Ленстры (эллиптические кривые)} & L_p[\tfrac12, \sqrt 2] \text{ (лекция~10)} \\
\text{решето числового поля} & L_N[\tfrac13, (64/9)^{1/3}], \ (64/9)^{1/3} \approx 1{,}923
\end{array}
\]
Решето числового поля (NFS) асимптотически быстрее, но сложнее и выигрывает лишь начиная с
$N$ порядка $10^{100}$"--~$10^{120}$; до этого рубежа лучший универсальный алгоритм "--- QS
(в варианте SIQS). Рекорды: в 1994~г. с помощью MPQS было разложено число RSA-129 (129 десятичных
знаков, около 600 добровольцев и 8 месяцев вычислений); рекорды последних лет установлены NFS:
RSA-240 (795 бит, 2019) и RSA-250 (829 бит, 2020)~\cite{boudot-240}. Поэтому модули RSA
короче $1024$ бит считаются небезопасными, а рекомендуемая длина "--- не меньше $2048$ бит.

\checkyourself{Повторите вывод теоремы~\ref{l6:qs-complexity} для метода Диксона с быстрой
проверкой гладкости. Какое $\beta$ получится и почему ответ хуже, чем для QS?}

\checkyourself{Почему решето заметно выгоднее пробного деления, хотя и то и другое находит
одни и те же гладкие значения? Во сколько раз (по порядку) различаются их затраты при
$B = L(N)^{1/2}$?}

\checkyourselfsection

\begin{furtherreading}
  \readitem[разд.~3.6--3.7]{hps}{разность квадратов, трёхшаговая процедура, гладкие числа,
    квадратичное решето и вывод его сложности; основной источник лекции}
  \readitem[разд.~5.1, гл.~6]{crandall-pomerance}{метод Ферма и метод Лемана; квадратичное
    решето, его варианты (MPQS, большие простые), линейная алгебра, решето числового поля}
  \readitem[разд.~4.5.4]{knuth-taocp2}{метод Ферма и его реализация, факторизация с помощью
    сравнений квадратов}
  \readitem[гл.~15]{galbraith}{гладкие числа, метод случайных квадратов (Диксона), оценки
    сложности в $L$-нотации}
  \readitem{pomerance-two-sieves}{история квадратичного решета и решета числового поля от их
    автора (статья в открытом доступе)}
  \readitem{pomerance-qs}{оригинальная статья о квадратичном решете}
  \readitem{dixon}{первый алгоритм с доказанной субэкспоненциальной оценкой}
  \readitem{cp-algorithms}{раздел «Integer factorization»: пробное деление, метод Ферма,
    $(p-1)$- и $\rho$-методы Полларда с кодом}
\end{furtherreading}
\end{lecture}
